%% =====================================================================
%% On the Value of Information in Monotone POMDPs
%% Canan Ulu and James E. Smith
%%
%% Working draft in the INFORMS OPRE author template (informs4.cls).
%% Build:  latexmk -pdf main.tex   (run from paper/; outputs to build/
%%         via latexmkrc if configured, or pass -output-directory=build)
%%
%% Draft conventions:
%%   \CU{...} purple  - comments by Canan Ulu (carried from notes)
%%   \JS{...} blue    - comments by Jim Smith (carried from notes)
%%   \Claude{...} teal - drafting notes by Claude: items needing the
%%                       authors' verification. Remove before submission.
%% =====================================================================
\documentclass[opre]{informs4}

\usepackage{eqndefns-left}
% Draft setting: relax the equation-width check from the published
% two-column width (240pt) to the manuscript width (470pt). Before
% submission, remove the two lines below and break any equations the
% checker then flags. (Tracked in notes/review-log.md.)
\setlength{\maxdseq}{470pt}
\EquationErrorBoxdimen=470pt
\RequirePackage{tgtermes}
\RequirePackage{newtxtext}
\RequirePackage{newtxmath}
\RequirePackage{bm}

\OneAndAHalfSpacedXII

\usepackage{array}
\usepackage{booktabs}
\usepackage{multirow}
\usepackage{pdflscape}
\usepackage{color}
\usepackage{tikz}
\usetikzlibrary{shapes.geometric}

% Natbib setup for author-year style (template default)
\usepackage{natbib}
 \bibpunct[, ]{(}{)}{,}{a}{}{,}%
 \def\bibfont{\small}%
 \def\bibsep{\smallskipamount}%
 \def\bibhang{24pt}%
 \def\newblock{\ }%
 \def\BIBand{and}%

\EquationsNumberedThrough
\TheoremsNumberedThrough
\ECRepeatTheorems

% ---------------------------------------------------------------------
% Notation macros (carried over from the authors' notes)
% ---------------------------------------------------------------------
\newcommand{\ev}[1]{\hspace{+0.20ex}\mathbb{E}\hspace{-0.40ex}\left[ \, #1 \, \right]}
\newcommand{\evpi}[1]{\hspace{.2em}\mathbb{E}_\pi\hspace{-.3em}\left[#1\right]}
\newcommand{\evsub}[2]{\hspace{.2em}\mathbb{E}_{#1}\hspace{-.3em}\left[#2\right]}
\newcommand{\iX}[0]{{\mathbb{X}}}
\newcommand{\iY}[0]{{\mathbb{Y}}}
\newcommand{\iXY}[0]{{\mathbb{XY}}}
\newcommand{\cev}[2]{\hspace{+0.20ex}\mathbb{E}\!\left[ \, #1 \, \middle| \, #2 \, \right]}
\newcommand\dif{\mathop{}\!\mathrm{d}}
\newcommand{\LR}{\succeq_{LR}}
\newcommand{\LRI}{\succeq_{LRI}}
\newcommand{\FOSD}{\succeq_{FOSD}}
\newcommand{\MQG}{\succeq_{MQG}}
\newcommand{\BW}{\succeq_{B}}
\newcommand{\Leh}{\succeq_{L}}

% ---------------------------------------------------------------------
% TikZ decision-tree styles (technology adoption figure)
% ---------------------------------------------------------------------
\tikzset{
  chance/.style={circle, draw, minimum size=4.5mm, inner sep=0pt},
  decision/.style={rectangle, draw, minimum size=4mm, inner sep=0pt},
  terminal/.style={isosceles triangle, isosceles triangle apex angle=55,
                   shape border rotate=180, draw, minimum size=2.6mm, inner sep=0pt},
  blab/.style={font=\large, anchor=south west, inner sep=1pt},
  nval/.style={font=\large, anchor=north, inner sep=2pt},
  leader/.style={densely dotted, thin},
  opt/.style={line width=1.1pt},
  plain/.style={},
}
\newcommand{\elbow}[4]{%
  \draw[#1] (#2.east) -- ([xshift=4.5mm]#2.east |- #3) -- (#3);
  \node[blab] at ([xshift=5mm]#2.east |- #3) {#4};}
\newcommand{\sigblock}[9]{%
  \pgfmathsetmacro{\yl}{#2}\pgfmathsetmacro{\yh}{#2-0.85}%
  \pgfmathsetmacro{\yr}{#2-1.7}\pgfmathsetmacro{\yc}{#2-0.425}%
  \pgfmathsetmacro{\yd}{#2-1.0625}%
  \node[decision] (d) at (2.7,\yd) {};
  \node[chance]   (c) at (5.2,\yc) {};
  \node[terminal] (tl) at (7.6,\yl) {};
  \node[terminal] (th) at (7.6,\yh) {};
  \node[terminal] (tr) at (5.2,\yr) {};
  \elbow{plain}{#1}{d}{#3}
  \elbow{#8}{d}{c}{Adopt}
  \elbow{#9}{d}{tr}{Reject}
  \elbow{plain}{c}{tl}{#4 Low}
  \elbow{plain}{c}{th}{#5 High}
  \node[nval, anchor=north east] at (d.south) {#6};
  \node[nval, anchor=north east] at (c.south) {#7};
  \draw[leader] (tl.east) -- (9.0,\yl);  \node[font=\scriptsize, anchor=east] at (9.6,\yl) {$-0.2$};
  \draw[leader] (th.east) -- (9.0,\yh);  \node[font=\scriptsize, anchor=east] at (9.6,\yh) {$0.4$};
  \draw[leader] (tr.east) -- (9.0,\yr);  \node[font=\scriptsize, anchor=east] at (9.6,\yr) {$0$};}
\newcommand{\payoffheader}[1]{%
  \node[font=\large\bfseries, anchor=east] at (9.6,#1) {Payoff};
  \draw[thin] (8.7,#1-0.22) -- (9.6,#1-0.22);}

% ---------------------------------------------------------------------
% Author margin-comment macros (removed before submission)
% ---------------------------------------------------------------------
\def\CU#1{\textcolor{purple}{#1}}
\def\JS#1{\textcolor{blue}{#1}}
\def\Claude#1{{\color[rgb]{0,0.5,0.5}[Claude: #1]}}
\renewcommand{\labelenumi}{\normalfont(\roman{enumi})}

%%%%%%%%%%%%%%%%
\begin{document}
%%%%%%%%%%%%%%%%

\RUNAUTHOR{Ulu and Smith}
\RUNTITLE{Value of Information in Monotone POMDPs}
\TITLE{On the Value of Information in Monotone POMDPs}

\ARTICLEAUTHORS{%
\AUTHOR{Canan Ulu}
\AFF{McDonough School of Business, Georgetown University, \EMAIL{cu50@georgetown.edu}}
\AUTHOR{James E. Smith}
\AFF{\Claude{Affiliation and email to confirm}}
}

\ABSTRACT{%
Partially observed Markov decision processes (POMDPs) in operations
research frequently possess monotone structure: value functions increase
and optimal actions shift upward as beliefs about the underlying state
improve. We study when one signal process is more valuable than another
in such monotone POMDPs. The classical answer, based on Blackwell's
comparison of experiments, requires one signal process to be a
state-independent garbling of the other---a demanding condition that
ranks few signal processes. Focusing on monotone POMDPs allows us to
relax this requirement. We say a signal process is \emph{LR-better} than
another if the latter can be obtained from the former by a
state-dependent garbling whose distributions are ordered by likelihood
ratios: the garbling in a lower state likelihood-ratio dominates the
garbling in a higher state. This order is implied by Blackwell
sufficiency and, in turn, implies the monotone quasi-garbling order of
\citet{kim2023comparing}, which uses first-order stochastic dominance
and characterizes the value of information in static monotone decision
problems. Likelihood-ratio dominance, unlike first-order stochastic
dominance, survives Bayesian updating and state transitions, which is
what makes the stronger order the appropriate one for dynamic problems:
we show that in a monotone POMDP, an LR-better signal process yields a
higher optimal value function at every horizon. We illustrate the
results with technology adoption, machine replacement, and medical
decision-making examples. \Claude{Abstract drafted; to be revised with
the paper.}
}%

\KEYWORDS{partially observed Markov decision processes; value of
information; Blackwell comparison of experiments; likelihood-ratio
dominance; monotone policies; garbling}

\maketitle

% ---------------------------------------------------------------------
% Section 1: Introduction.
% Drafted by Claude following the paragraph plan in Outline.docx.
% Literature items the outline asked Claude to research are flagged
% with \Claude{} comments; candidates with summaries are in
% notes/literature-candidates.md for the authors to select from.
% ---------------------------------------------------------------------
\section{Introduction}\label{sec:intro}

Partially observed Markov decision processes (POMDPs) are pervasive in
operations research \citep{krishnamurthy2016partially}. In a POMDP, a
decision maker (DM) repeatedly chooses actions under uncertainty about
an underlying state, observing signals that partially reveal the state
and updating her beliefs over time. Examples include technology
adoption, where a consumer or firm decides whether to adopt a new
technology of uncertain benefit or wait and gather more information
\citep{mccardle1985information, lippman1987does, ulu2009uncertainty};
machine maintenance, where a DM decides when to inspect or replace
equipment whose condition is only partially observed
\citep{white1979optimal, maillart2006maintenance}; inventory control
with unobserved demand states \citep{treharne2002adaptive}; and
healthcare applications such as prostate biopsy referral decisions
\citep{zhang2012optimization} and patient-type adaptive treatment
planning \citep{skandari2021patient}. See
\citet{krishnamurthy2016partially} for many other examples.

Many POMDP models of interest have \emph{monotone} structure: the
optimal action increases as the DM's beliefs increase in a stochastic
sense. In technology adoption models, as beliefs about the benefits of
the technology improve, optimal policies move toward adoption
\citep{mccardle1985information, ulu2009uncertainty}. In machine
maintenance models, as beliefs about the state of the machine worsen,
optimal policies move toward replacement \citep{white1979optimal,
lovejoy1987some}. And as beliefs point toward cancer, optimal policies
move toward a biopsy \citep{zhang2012optimization}. A substantial
literature provides sufficient conditions for such structure:
\citet{lovejoy1987some} gives conditions for the optimal value
function of a finite POMDP to be monotone in the prior, with priors
ordered by likelihood ratios, for the optimal policy in a machine
replacement problem to be monotone, and, in the general case, for the
optimal policy to be bounded below by a monotone function.
\Claude{Your outline asks for additional structural-results
references, especially by Vikram Krishnamurthy. A candidate list with
summaries is in notes/literature-candidates.md for you to select
from; once you choose, I will fold them into this paragraph.}

How should one compare information sources in a POMDP? The classical
tool is Blackwell's comparison of experiments
\citep{blackwell1951comparison}: if, for each action, the signal
process is garbled by adding noise that is independent of the state,
the resulting POMDP has a lower optimal value than the original.
Establishing this value comparison uses the convexity of the POMDP
value function in the prior \citep{smallwood1973optimal} and requires
no monotonicity of the value function or optimal policies. Blackwell's
order is powerful precisely because it is preference free---a garbling
is worse in \emph{every} decision problem---but this universality
makes it a very strong requirement, and few pairs of information
structures can be ranked by it \citep{lehmann1988comparing,
kim2023comparing}. A large literature has therefore sought weaker
orders that rank more information structures by restricting the class
of decision problems, most prominently to \emph{monotone} decision
problems: \citet{lehmann1988comparing} introduces the accuracy order
for single-decision problems with monotone likelihood-ratio signals;
\citet{persico2000information} and \citet{quah2009comparative} extend
the reach of Lehmann's comparison to broader classes of monotone
payoffs; and \citet{jewitt2007information} connects Lehmann's order to
Blackwell dominance on dichotomies. \citet{athey2018value} study
monotone decision problems in a related but different setup, where
information structures are compared for a \emph{fixed} prior belief:
their orderings are prior dependent, which permits finer comparisons
when the prior is known, whereas the orders we study---like Blackwell's
and Lehmann's---rank signal processes uniformly across priors.
\Claude{Two items to verify here: (i) the outline asks for citations
supporting the ``Blackwell is too strong'' criticism---I cite Lehmann
and Kim, both of whom make the point explicitly; add others if you
have favorites. (ii) The outline asks whether Krishnamurthy has papers
using Lehmann's order: I could not identify any---his structural
POMDP results use Blackwell dominance of observation kernels (see
notes/literature-candidates.md)---please confirm before we assert
anything in text.}

All of these weaker orders, however, were developed for static,
single-decision problems. In this paper, we ask the corresponding
question for dynamic problems: when is one signal process worth more
than another in a \emph{monotone POMDP}? Because we focus on a smaller
class of POMDPs, we can obtain a weaker order among information
structures that is implied by Blackwell sufficiency. We say a signal
process $\iX$ is \emph{LR-better} than a signal process $\iY$ if $\iY$
can be obtained from $\iX$ by adding a state-dependent garbling such
that the garbling distribution at a lower state likelihood-ratio (LR)
dominates the garbling distribution at a higher state---that is, the
added noise systematically pushes signals upward in low states and
downward in high states, degrading the signal's ability to
discriminate between states. Our order is a likelihood-ratio
strengthening of the \emph{monotone quasi-garbling} order of
\citet{kim2023comparing}, who studies monotone decision problems with
a single decision and shows that monotone quasi-garbling---where the
garbling distribution at a lower state first-order stochastically
dominates the garbling at a higher state---is a necessary and
sufficient condition for all decision makers with monotone decision
problems to obtain higher expected payoffs. LR dominance implies
first-order stochastic dominance, so the LR-better order implies
monotone quasi-garbling. The strengthening is exactly what the dynamic
setting demands: LR dominance survives Bayesian updating and state
transitions in a POMDP, whereas first-order stochastic dominance does
not \citep[see][for a discussion]{ulu2009uncertainty}.

Our main result (Theorem~\ref{thm:main}) shows that in a monotone
POMDP---one with monotone rewards, monotone likelihood-ratio signal
and transition processes, and supermodular structure delivering
monotone optimal policies---an LR-better signal process yields a
higher optimal value function at every horizon. We also map the
relationships among the orders (Section~\ref{sec:orders}): under the
monotone likelihood ratio property, Blackwell's order implies the
LR-better order, which implies monotone quasi-garbling, which is in
turn equivalent to Lehmann's accuracy order. We illustrate the
framework with technology adoption, machine replacement, and medical
decision-making models, present Normal, Poisson, and Uniform families
of examples in which LR-better comparisons arise naturally, and show
by example that when the state dependence of the noise runs in the
opposite direction---signals from better states are \emph{more} likely
to be observed---the two signal processes may be genuinely
incomparable, with value functions that cross as a function of the
prior.

The rest of the paper is organized as follows.
Section~\ref{sec:model} presents the monotone POMDP model and its
structural properties. Section~\ref{sec:comparison} introduces the
LR-better order and establishes the main value-comparison theorem.
Section~\ref{sec:orders} relates the LR-better order to the Blackwell,
Lehmann, and monotone quasi-garbling orders.
Section~\ref{sec:applications} presents applications and examples,
and Section~\ref{sec:conclusion} concludes.

% ---------------------------------------------------------------------
% Section 2: Monotone POMDPs
% Source: notes/source-material/LRMQG.tex, Sections "Setup", "Beliefs",
% "The Dynamic Program"; guiding text drafted by Claude.
% Propositions stated without proof per the authors' instruction;
% surrounding text points the reader to the relevant literature.
% ---------------------------------------------------------------------
\section{Monotone POMDPs}\label{sec:model}

\subsection{States, Actions, and Beliefs}\label{sec:setup}

We consider a discrete-time, finite-horizon POMDP. The state of the
world is denoted $\theta \in \Theta$, where $\Theta \subseteq \Re$ is a
fully ordered set; for example, $\theta$ may represent the benefit of a
new technology in a technology adoption model or the condition of a
machine in a machine replacement model. In each period, the decision
maker (DM) chooses an action $a$ from a finite and fully ordered set
$A = \{a_0, a_1, \ldots, a_N\}$, where $a_0 \leq a_1 \leq \cdots \leq
a_N$. The DM cannot observe the state directly. Her beliefs about
$\theta$ are described by a probability distribution; for ease of
notation, we assume that this distribution is continuous with a density
$\pi$ over $\Theta$. For discrete state spaces, we can interpret $\pi$
as a probability mass function and replace integrals with sums
\citep{ulu2009uncertainty}.

We index periods by the number of periods remaining, $k$. The state may
change over time: if the state in period $k$ is $\theta_k$ and the DM
takes action $a$, the next-period state $\theta_{k-1}$ is drawn from the
conditional distribution $g(\theta_{k-1}|\theta_k, a)$. Given a prior
$\pi$ on the period-$k$ state, the DM's \emph{next-period prior} is
\begin{equation}
\eta(\theta_{k-1}; \pi, a)
  = \int_{\theta_k \in \Theta} g(\theta_{k-1}|\theta_k, a)\,
    \pi(\theta_k) \dif \theta_k \ . \label{eq:eta}
\end{equation}
Information arrives through a \emph{signal process} $\iX$ that generates
signals $x \in X$, where $X$ is fully ordered, with likelihood function
$L_\iX(x|\theta, a)$ describing the probability of observing signal $x$
given (next-period) state $\theta$ and action $a$. The \emph{predictive
distribution} for signals and the \emph{posterior distribution} on the
state given a signal are then
\begin{align}
f_\iX(x;\pi, a) &= \int_\theta L_\iX(x|\theta,a)\, \eta(\theta; \pi, a)
  \dif \theta \ , \label{eq:predictive}\\
\Pi_\iX(\theta; \pi, x, a) &=
  \frac{L_\iX(x|\theta,a)\, \eta(\theta; \pi, a)}{f_\iX(x;\pi, a)} \ .
  \label{eq:posterior}
\end{align}
Thus, within a period, beliefs evolve from the prior $\pi$ to the
next-period prior $\eta(\,\cdot\,; \pi, a)$ through the state
transition, and then to the posterior $\Pi_\iX(\,\cdot\,; \pi, x, a)$
after the signal $x$ is observed; this posterior serves as the prior for
the next period. Because the distributions themselves will serve as
arguments of value functions, we will frequently suppress the state
argument and write, e.g., $\eta(\pi,a)$, $f_\iX(\pi,a)$, and
$\Pi_\iX(\pi,x,a)$ when we consider these distributions as functions of
the prior $\pi$, the signal $x$, and the action $a$.

When studying monotonicity properties involving distributions, we focus
on likelihood-ratio (LR) dominance as a stochastic order.

\begin{definition}[LR order and MLR property]\label{def:LR}\quad
\begin{enumerate}
\item A distribution $\pi_2$ \emph{LR-dominates} $\pi_1$
  ($\pi_2 \LR \pi_1$) if for all $\theta_2 \geq \theta_1$,
  \begin{equation}
  \frac{\pi_2(\theta_2)}{\pi_1(\theta_2)} \geq
  \frac{\pi_2(\theta_1)}{\pi_1(\theta_1)} \ . \label{eq:LRdef}
  \end{equation}
\item $L_\iX(x|\theta, a)$ satisfies the \emph{MLR property} if
  $L_\iX(\,\cdot\,|\theta_2, a) \LR L_\iX(\,\cdot\,|\theta_1, a)$
  whenever $\theta_2 \geq \theta_1$. Similarly, $g(\theta_{k-1}|
  \theta_k, a)$ satisfies the MLR property if
  $g(\,\cdot\,|\theta_k^2, a) \LR g(\,\cdot\,|\theta_k^1, a)$ whenever
  $\theta_k^2 \geq \theta_k^1$.
\end{enumerate}
\end{definition}

The ratios in \eqref{eq:LRdef} may be interpreted as being infinite when
the denominators are zero; equivalently, $\pi_2 \LR \pi_1$ if
$\pi_2(\theta_2)\pi_1(\theta_1) \geq \pi_2(\theta_1)\pi_1(\theta_2)$
for all $\theta_2 \geq \theta_1$. LR dominance implies first-order
stochastic dominance (FOSD) and, unlike FOSD, LR dominance survives
Bayesian updating and the state transition process, provided the
likelihoods and state transitions satisfy the MLR property; see
\citet{ulu2009uncertainty} for a discussion of this distinction and
examples where FOSD fails to be preserved. The following proposition
records these properties in our setting. The results are well known
(see, e.g., \citealt{whitt1979note}; \citealt{lovejoy1987some},
Lemma~1.2, for the finite-state case; and \citealt{ulu2009uncertainty},
Propositions~3.2 and 3.4 and \S8), and we state them without proof.

\begin{proposition}[Monotonicity of beliefs]\label{prop:monotone_beliefs}
For any action $a$, if $L_\iX(x|\theta,a)$ and
$g(\theta_{k-1}|\theta_k,a)$ satisfy the MLR property, then:
\begin{enumerate}
\item $\pi_2 \LR \pi_1$ implies $\eta(\pi_2,a) \LR \eta(\pi_1,a)$,
  $f_\iX(\pi_2,a) \LR f_\iX(\pi_1,a)$, and
  $\Pi_\iX(\pi_2, x, a) \LR \Pi_\iX(\pi_1, x, a)$ for all $x$.
\item For any $\pi$, $x_2 \geq x_1$ if and only if
  $\Pi_\iX(\pi,x_2,a) \LR \Pi_\iX(\pi,x_1,a)$.
\end{enumerate}
\end{proposition}

Part (i) says that more optimistic priors lead to more optimistic
next-period priors, signal distributions, and posteriors; part (ii) says
that more favorable signals lead to more optimistic posteriors. This
``survival'' of the LR order through Bayesian updating and state
transitions is the key property that drives the recursive arguments in
the dynamic program below.
\CU{Mention more properties of the order? e.g., single-crossing
preservation/dominance??}

\subsection{The Dynamic Program}\label{sec:dp}

The DM earns a reward $r(\theta, a)$ when taking action $a$ in state
$\theta$; given beliefs $\pi$, the expected reward is
\begin{equation}
r(\pi,a) = \int_{\theta \in \Theta} r(\theta,a)\, \pi(\theta)
  \dif \theta \ . \label{eq:expreward}
\end{equation}
Rewards are discounted with a discount factor $\delta \in [0,1]$. The
DM's optimal value function with $k$ periods remaining is defined by the
dynamic programming recursion: let $V_{\iX,0}^*(\pi) = 0$ and, for
$k > 0$,
\begin{align}
V_{\iX,k}(\pi,a) &= r(\pi,a) +
  \delta \ev{V_{\iX,k-1}^*(\Pi_\iX(\pi,\tilde{x},a))} \ ,
  \label{eq:Vka}\\
V_{\iX,k}^*(\pi) &= \max_{a\in A} \left\{ V_{\iX,k}(\pi,a) \right\} \ ,
  \label{eq:Vk}
\end{align}
where the expected continuation value is taken with respect to the
predictive distribution,
\begin{equation}
\ev{V_{\iX,k-1}^*(\Pi_\iX(\pi,\tilde{x},a))}
  = \int_{x\in X} V_{\iX,k-1}^*(\Pi_\iX(\pi,x,a))\,
    f_\iX(x; \pi, a) \dif x \ . \label{eq:contval}
\end{equation}
An \emph{action selection} $\alpha_{k}(\pi)$ chooses an action
$a \in A$ in period $k$ for a given prior $\pi$; a \emph{policy}
$\boldsymbol{\alpha} = (\alpha_{k}, \ldots, \alpha_{0})$ is a sequence
of action selections. An action selection $\alpha_{k}^*$ is optimal if
$V_{\iX,k}^*(\pi) = V_{\iX,k}(\pi,\alpha_{k}^*(\pi))$ for all $\pi$,
and a policy $\boldsymbol{\alpha}^*$ is optimal if each of its action
selections is optimal.

We say a function $u(\pi)$ defined on distributions is
\emph{LR-increasing} if $u(\pi_2) \geq u(\pi_1)$ whenever
$\pi_2 \LR \pi_1$. Our first structural result states that, with
increasing rewards and MLR likelihoods and transitions, the value
functions are LR-increasing: LR improvements in the DM's beliefs make
the DM better off. The proof is by induction on $k$, mirroring the
proof of Proposition~1 of \citet{lovejoy1987some} in the finite-state
case and Proposition~3.6 of \citet{ulu2009uncertainty} in the technology
adoption model: expected rewards \eqref{eq:expreward} are LR-increasing
because $r(\theta,a)$ is increasing in $\theta$ and LR dominance implies
FOSD; the expected continuation value \eqref{eq:contval} is
LR-increasing by Proposition~\ref{prop:monotone_beliefs} and the
induction hypothesis; and the maximum \eqref{eq:Vk} of LR-increasing
functions is LR-increasing. We omit the details.

\begin{proposition}[Monotonicity of value functions]\label{prop:monotone_values}
If, for all actions $a$, $r(\theta, a)$ is increasing in $\theta$ and
$L_\iX(x|\theta,a)$ and $g(\theta_{k-1}|\theta_k,a)$ satisfy the MLR
property, then:
\begin{enumerate}
\item $r(\pi,a)$ and $V_{\iX,k}(\pi,a)$ are LR-increasing for all $a$
  and $k$, and
\item $V_{\iX,k}^*(\pi)$ is LR-increasing for all $k$.
\end{enumerate}
\end{proposition}

\begin{remark}[Comparison with \citealt{lovejoy1987some}]\label{rem:lovejoy}
Proposition~\ref{prop:monotone_values} is a general-state-space analog
of Proposition~1 of \citet{lovejoy1987some}, and the assumptions map
to his as follows. We assume that the state transition process
$g(\theta_{k-1}|\theta_k,a)$ satisfies the MLR property; Lovejoy's
Lemma~1.3(1) shows that this (his TP$_2$ condition on the transition
matrix) is a sufficient condition for the monotonicity of next-period
priors required in his Lemma~1.2(2)---a condition that is, in his
finite-state setting, also necessary. Lovejoy's Lemma~1.3(2) provides
conditions under which next-period priors are LR-increasing in the
action; neither Lovejoy's Proposition~1 nor our
Proposition~\ref{prop:monotone_values} requires this assumption.
\end{remark}

\subsection{Monotone Optimal Policies}\label{sec:policies}

To establish the monotonicity of optimal policies, we introduce the
notion of increasing differences, also known as supermodularity.

\begin{definition}[Increasing differences]\label{def:incdiff}\quad
\begin{enumerate}
\item A function $u(\theta,a)$ has \emph{increasing differences} if
  $u(\theta_2,a_2)-u(\theta_2,a_1) \geq u(\theta_1,a_2)-u(\theta_1,a_1)$
  whenever $\theta_2 \geq \theta_1$ and $a_2 \geq a_1$.
\item A function $u(\pi,a)$ has \emph{LR-increasing differences} if
  $u(\pi_2,a_2)-u(\pi_2,a_1) \geq u(\pi_1,a_2)-u(\pi_1,a_1)$ whenever
  $\pi_2 \LR \pi_1$ and $a_2 \geq a_1$.
\item The state transitions defined by $L_\iX(x|\theta,a)$ and
  $g(\theta_{k-1}|\theta_k,a)$ satisfy a \emph{stochastic version of
  increasing differences} if
  \begin{equation}
  \ev{u(\Pi_\iX(\pi_2,\tilde{x},a_2))}-\ev{u(\Pi_\iX(\pi_2,\tilde{x},a_1))}
  \geq
  \ev{u(\Pi_\iX(\pi_1,\tilde{x},a_2))}-\ev{u(\Pi_\iX(\pi_1,\tilde{x},a_1))}
  \label{eq:stochincdiff}
  \end{equation}
  for all LR-increasing functions $u(\pi)$, whenever $\pi_2 \LR \pi_1$
  and $a_2 \geq a_1$.
\end{enumerate}
\end{definition}

Note that the functions $u(\pi)$ in part (iii) are assumed to be
LR-increasing; this could be replaced by other conditions, e.g.,
LR-increasing and convex, or many other choices.

The next lemma records two standard consequences of increasing
differences. Part (i) shows that increasing differences in the state
carry over to LR-increasing differences in beliefs when taking
expectations. Part (ii) is the standard monotone comparative statics
result---increasing differences imply increasing optimal
selections---adapted to the LR order on beliefs; it is a special case
of Lemma~6.1 of \citet{topkis1978minimizing}, and parallels Lemma~2.1
of \citet{lovejoy1987some}. \Claude{The statement of part (ii) below
was drafted by me---the notes contained only the placeholder
``increasing differences to increasing policies.'' Please verify the
statement (in particular, the use of the greatest maximizer to obtain a
well-defined monotone selection).}

\begin{lemma}\label{lemma:incdiff}\quad
\begin{enumerate}
\item If $u(\theta,a)$ has increasing differences, then
  $u(\pi,a) = \int_{\theta \in \Theta} u(\theta,a) \pi(\theta)
  \dif \theta$ has LR-increasing differences.
\item If $u(\pi,a)$ has LR-increasing differences, then the greatest
  maximizer $\alpha^u(\pi) = \max\, \mathrm{argmax}_{a \in A}\,
  u(\pi,a)$ is LR-increasing: $\pi_2 \LR \pi_1$ implies
  $\alpha^u(\pi_2) \geq \alpha^u(\pi_1)$. In particular, if
  $L_\iX(x|\theta,a)$ and $g(\theta_{k-1}|\theta_k,a)$ satisfy the MLR
  property, then for any prior $\pi$ and action $a$, the induced signal
  selection $\alpha^u(\Pi_\iX(\pi,x,a))$ is increasing in the observed
  signal $x$.
\end{enumerate}
\end{lemma}

The second claim in part (ii) follows by combining the first claim with
Proposition~\ref{prop:monotone_beliefs}(ii): more favorable signals
lead to LR-larger posteriors and hence to larger selected actions. We
can now state the monotone policy result. As with
Proposition~\ref{prop:monotone_values}, the proof of part (i) is by
induction on $k$---the reward term has LR-increasing differences by
Lemma~\ref{lemma:incdiff}(i) and the continuation term satisfies
\eqref{eq:stochincdiff} with $u = V_{\iX,k-1}^*$, which is
LR-increasing by Proposition~\ref{prop:monotone_values}---and part (ii)
follows from part (i) via Lemma~\ref{lemma:incdiff}(ii). We omit the
details. \Claude{Confirm you want this proof sketch in text rather
than a full proof in an appendix; OR will ultimately expect complete
proofs in an e-companion.}

\begin{proposition}[Monotonicity of optimal policies]\label{prop:monotone_policies}
Suppose (i) $r(\theta,a)$ is increasing in $\theta$ for all $a$ and has
increasing differences, and (ii) $L_\iX(x|\theta,a)$ and
$g(\theta_{k-1}|\theta_k,a)$ satisfy the MLR property for all $a$ and
satisfy the stochastic version of increasing differences. Then:
\begin{enumerate}
\item $V_{\iX,k}(\pi,a)$ has LR-increasing differences for all $k$.
\item There exists an optimal policy $\boldsymbol{\alpha}_{\iX}^*$ that
  is LR-increasing.
\end{enumerate}
\end{proposition}

Propositions \ref{prop:monotone_values} and
\ref{prop:monotone_policies} describe the class of models we will call
\emph{monotone POMDPs}: models in which values increase and optimal
actions shift upward as beliefs improve in the LR order. As discussed
in the introduction, many POMDP models in operations research fall in
this class. In the remainder of the paper, we study when one signal
process is more valuable than another in such models.

% ---------------------------------------------------------------------
% Section 3: Comparing signal processes; LR-better order; main theorem.
% Source: LRMQG.tex "Comparing Signal Processes" (definition, lemma,
% and the authors' proof of the lemma). The main theorem statement and
% its induction proof were drafted by Claude per the authors'
% instruction and are flagged for verification.
% ---------------------------------------------------------------------
\section{Comparing Signal Processes}\label{sec:comparison}

We now turn to the central question of the paper: when is one signal
process more valuable than another in a monotone POMDP? Consider two
signal processes $\iX$ and $\iY$ generating signals $x \in X$ and
$y \in Y$ with likelihoods $L_\iX(x|\theta,a)$ and $L_\iY(y|\theta,a)$,
respectively, where $X$ and $Y$ are fully ordered. The two processes
provide information about the same underlying states, and the two
POMDPs are otherwise identical: they share the same state space,
actions, rewards, state transitions $g$, and discount factor, and the
value functions $V_{\iY,k}(\pi,a)$ and $V_{\iY,k}^*(\pi)$ are defined
by the recursion \eqref{eq:Vka}--\eqref{eq:Vk} with $\iY$ in place of
$\iX$. If the DM observes signals from $\iX$ instead of $\iY$, when is
she better off?

The classical answer is given by Blackwell's comparison of experiments
\citep{blackwell1951comparison}: if $\iY$ can be obtained from $\iX$ by
adding noise that is independent of the state---that is, if $\iY$ is a
\emph{garbling} of $\iX$---then $\iX$ is more valuable in \emph{every}
decision problem. In the POMDP setting, the corresponding result---that
garbling the signal process for each action lowers the optimal value
function---follows from the convexity of the value function in the
prior \citep{smallwood1973optimal} and requires no monotonicity
structure. The price of this generality is that the Blackwell order is
very incomplete: most pairs of signal processes cannot be ranked. By
restricting attention to monotone POMDPs, we can obtain a weaker
order---one that ranks more signal processes---that is implied by
Blackwell sufficiency.

\subsection{LR-Better Signal Processes}\label{sec:lr-better}

\begin{definition}[LR-better signal processes]\label{def:LRI}
We say that a signal process $\iX$ is \emph{LR-better} than signal
process $\iY$ (denoted $\iX \LRI \iY$) if there exists a conditional
probability distribution $P(y|x,\theta)$ such that
\begin{enumerate}
\item $L_\iY(y|\theta) = \int P(y|x,\theta) L_\iX(x|\theta) \dif x$
  for all $y \in Y$ and $\theta$, and
\item $P(\,\cdot\,|x,\theta_1) \LR P(\,\cdot\,|x,\theta_2)$ for all
  $x \in X$ and $\theta_2 \geq \theta_1$.
\end{enumerate}
\end{definition}

\Claude{In the dynamic model the likelihoods depend on the action $a$.
I have stated the definition with $a$ suppressed, as in your notes; I
assume the intended reading is that conditions (i)--(ii) hold for each
action $a$, with a garbling kernel $P_a(y|x,\theta)$ that may depend on
$a$. Please confirm, and I will make the dependence explicit if you
prefer.}

Condition (i) says that $\iY$ is obtained from $\iX$ by adding noise
through the garbling kernel $P$; unlike in Blackwell's comparison, the
noise may depend on the state $\theta$. Condition (ii) disciplines this
state dependence: the garbling distribution in a lower state
LR-dominates the garbling distribution in a higher state. In other
words, the added noise is \emph{reversely monotone}---it tends to
produce higher $y$ signals in lower states and lower $y$ signals in
higher states, and thus degrades the information carried by the
signal. When $P(y|x,\theta)$ does not depend on $\theta$, condition
(ii) holds trivially (every distribution LR-dominates itself), and
$\iX \LRI \iY$ reduces to Blackwell's garbling condition. We study the
relationship between the LR-better order and other information orders
in Section~\ref{sec:orders}.

\subsection{Main Result}\label{sec:main-result}

Our main result states that in a monotone POMDP, an LR-better signal
process yields a higher optimal value function at every horizon. The
key step is the following lemma, which compares the expected optimal
continuation values under the two signal processes for a common payoff
function with LR-increasing differences.

\begin{lemma}[LR-better continuation values]\label{lemma:LRI-ContVal}
Suppose
\begin{enumerate}
\item $u(\pi,a)$ has LR-increasing differences,
\item $L_\iX(x|\theta,a)$, $L_\iY(y|\theta,a)$, and
  $g(\theta_{k-1}|\theta_k,a)$ satisfy the MLR property for all $a$,
  and
\item $\iX \LRI \iY$,
\end{enumerate}
and let $u^*(\pi) = \max_{a\in A} u(\pi, a)$. Then, for all $\pi$ and
$a$,
\begin{equation}
\ev{u^*(\Pi_\iX(\pi,\tilde{x},a))} \geq
\ev{u^*(\Pi_\iY(\pi,\tilde{y},a))} \ . \label{eq:lemma_result}
\end{equation}
\end{lemma}

The idea of the proof is as follows. We construct a joint signal model
in which the $\iY$ signal is generated from the $\iX$ signal through
the garbling kernel $P$, so that observing $(x,y)$ is at least as
informative as observing $y$ alone. An optimal action selection based
on $y$ alone is feasible in the joint model, and---because the
selection is monotone and the garbled posteriors move \emph{down} in
the LR order as $y$ increases for a fixed $x$---we show that, given
$x$, a \emph{constant} action outperforms it. Averaging over signals
then delivers \eqref{eq:lemma_result}. The formal proof follows.

\begin{proof}{Proof.}
Fix $\pi$ and $a$ and let $\alpha_\iY^*(\Pi_\iY(\pi, y,a))$ be an
optimal action selection that is increasing in $y$ and such that
$$u^*(\Pi_\iY(\pi,\tilde{y},a)) =
u(\Pi_\iY(\pi,\tilde{y},a),\alpha_\iY^*(\Pi_\iY(\pi, \tilde{y},a)))\ .$$
Such an optimal increasing action selection exists because $u(\pi,a)$
is assumed to have LR-increasing differences and the posteriors
$\Pi_\iY(\pi,y,a)$ are LR-increasing in $y$; this follows from the
supermodularity argument of Lemma~\ref{lemma:incdiff}(ii). Similarly,
let $\alpha_\iX^*(\cdot)$ be an optimal action selection for $\iX$.

We consider a joint signal model $\iXY$, taking
$L_\iXY((x,y)|\theta) = P(y|x,\theta) L_\iX(x|\theta)$, with predictive
and posterior distributions $f_\iXY((x,y); \pi,a)$ and
$\Pi_\iXY(\theta; (x,y),\pi,a)$ defined as in
\eqref{eq:predictive}--\eqref{eq:posterior}. Note that $\iX \LRI \iY$
implies that $\Pi_\iXY(\theta; (x,y),\pi,a)$ is LR-decreasing in $y$
for fixed $x$, $\pi$, and $a$.

To ease the notation, we will suppress the fixed prior $\pi$ and the
action $a$, as these are constant throughout, taking
$\Pi_\iXY(x,y) = \Pi_\iXY(\pi,(x,y),a)$ and
$f_\iXY(x,y) = f_\iXY((x,y); \pi,a)$ to be the posterior and predictive
distributions for $\iXY$. Similarly, let
$\Pi_\iX(x) = \Pi_\iX(\pi,x,a)$ and $\Pi_\iY(y) = \Pi_\iY(\pi,y,a)$ be
the posteriors for $\iX$ and $\iY$, and denote the selection function
$\alpha_\iY^*(y) = \alpha_\iY^*(\Pi_\iY(\pi, y,a))$.

Since the action selection $\alpha_\iY^*$ ignores the $x$ signal,
\begin{align}
\ev{u^*(\Pi_\iY(\tilde{y}))}
&= \ev{u(\Pi_\iY(\tilde{y}),\alpha_\iY^*(\tilde{y}))} \notag\\
&= \ev{u(\Pi_\iXY(\tilde{x},\tilde{y}),\alpha_\iY^*(\tilde{y}))}
  \notag\\
&= \ev{\cev{u(\Pi_\iXY(\tilde{x},\tilde{y}),\alpha_\iY^*(\tilde{y}))}
  {\tilde{x}}} \ , \label{eq:nested_expectation}
\end{align}
where in the last expression we condition on a given signal $\tilde{x}$
in the inner expectation, and then take expectations over $\tilde{x}$.

We will focus on the inner expectation in
\eqref{eq:nested_expectation} with a given signal $x$ and show that,
given this $x$, there exists a constant action $a'$ (i.e., independent
of $y$) that outperforms $\alpha_\iY^*$. Because $\alpha_\iY^*$ is
increasing, there are thresholds $y_0 \leq y_1 \leq \ldots \leq
y_{N+1}$ such that action $a_i$ is chosen for signals in
$[y_i,y_{i+1})$ (where the interval is empty if $y_i = y_{i+1}$). Then
\begin{align}
&\cev{u(\Pi_\iXY(x,\tilde{y}),\alpha_\iY^*(\tilde{y}))}{x} \notag\\
&\qquad = \sum_{i=0}^N \int_{y_{i}}^{y_{i+1}} u(\Pi_\iXY(x,y),a_i)\,
  f_\iXY(x,y) \dif y \ . \label{eq:regions}
\end{align}
We will identify an improving constant action $a'$ by sequentially
considering the signal regions $[y_i,y_{i+1})$ for $i = 1, \ldots, N$
and improving the action selection as we go. First, take $i=1$ and
consider a candidate constant action $a' = a_0$, the action selected by
$\alpha_\iY^*$ in region $[y_0,y_1)$. Consider the sign of the
difference
\begin{align}
D_i = \int_{y_i}^{y_{i+1}}
  &\left(u(\Pi_\iXY(x,y),a_i) - u(\Pi_\iXY(x,y),a')\right) \notag\\
  &\quad \times f_\iXY(x,y) \dif y \ , \label{eq:Di}
\end{align}
where $a_i$ is the action selected by $\alpha_\iY^*$ in region $i$.
First, if $D_i \leq 0$, then selecting action $a'$ instead of $a_i$ in
region $i$ would lead to an improvement in the performance of the
action selection. Second, consider the case where $D_i \geq 0$.
Because $u(\pi,a)$ has LR-increasing differences and $\Pi_\iXY(x, y)$
is LR-decreasing in $y$,
$$u(\Pi_\iXY(x,y),a_i) - u(\Pi_\iXY(x,y),a')$$
is decreasing in $y$. If $D_i \geq 0$, this implies that choosing
action $a_i$ instead of action $a'$ over the interval $[y_0,y_i)$
would lead to an improvement in performance; in this case, we replace
$a'$ with $a_i$. Thus, considering both cases and selecting the
corresponding $a'$, we have an action selection that chooses $a'$ over
the regions $[y_0, y_{i+1})$ and improves on $\alpha_\iY^*$, for this
given $x$. Repeating this process for regions $i+1, \ldots, N$, we
arrive at a constant action $a'$ that implicitly depends on $x$ but is
independent of $y$ and improves on $\alpha_\iY^*$. Making the implicit
dependence on $x$ explicit, we write this action selection $a'(x)$ and
have
\begin{align}
\cev{u(\Pi_\iXY(x,\tilde{y}),\alpha_\iY^*(\tilde{y}))}{x}
&\leq \sum_{i=0}^N \int_{y=y_{i}}^{y_{i+1}}
  u(\Pi_\iXY(x,y),a'(x))\, f_\iXY(x,y) \dif y \notag\\
&= \cev{u(\Pi_\iXY(x,\tilde{y}),a'(x))}{x} \notag\\
&= u(\Pi_\iX(x),a'(x)) \ , \label{eq:constant_action}
\end{align}
where the final equality follows because, with the action fixed at
$a'(x)$, averaging the posteriors $\Pi_\iXY(x,\tilde y)$ over
$\tilde{y}$ given $x$ recovers the posterior $\Pi_\iX(x)$.
\Claude{I tightened the final step here: your notes wrote the last line
as $\ev{u(\Pi_\iX(x),a'(x)) \mid x}$; since $u$ is linear in the
belief argument (it is an expectation of a payoff under the belief),
averaging the $\iXY$ posteriors over $y$ given $x$ yields
$\Pi_\iX(x)$. This linearity deserves a check: if $u(\pi,a)$ is an
arbitrary function with LR-increasing differences (e.g., a value
function including continuation terms), it need not be linear in
$\pi$, and this step instead requires the tower property
$\ev{u(\Pi_\iXY(x,\tilde y), a')\,|\,x} = u(\Pi_\iX(x), a')$ to be
justified---please verify how you want this argued.}
Then taking expectations over $\tilde{x}$ and continuing from
\eqref{eq:nested_expectation}, we have
\begin{align}
\ev{u^*(\Pi_\iY(\tilde{y}))}
&= \ev{u(\Pi_\iY(\tilde{y}),\alpha_\iY^*(\tilde{y}))} \notag\\
&= \ev{\cev{u(\Pi_\iXY(\tilde{x},\tilde{y}),
  \alpha_\iY^*(\tilde{y}))}{\tilde{x}}} \notag\\
&\leq \ev{u(\Pi_\iX(\tilde{x}),a'(\tilde{x}))} \notag\\
&\leq \ev{u(\Pi_\iX(\tilde{x}),\alpha_\iX^*(\tilde{x}))} \notag\\
&= \ev{u^*(\Pi_\iX(\tilde{x}))} \ ,
\end{align}
where the final inequality follows from the fact that $a'(\tilde{x})$
is feasible but not necessarily an optimal action selection for $\iX$.
\Halmos
\end{proof}

With Lemma~\ref{lemma:LRI-ContVal} in hand, the value comparison
follows by induction on the horizon. \Claude{The statement and proof
of Theorem~\ref{thm:main} below are my draft, per your instruction;
please verify both. Note in particular that the stochastic
increasing-differences condition is imposed only on the $\iX$ process:
the induction applies Lemma~\ref{lemma:LRI-ContVal} with
$u = V_{\iX,k-1}$, whose LR-increasing differences come from
Proposition~\ref{prop:monotone_policies}(i) for $\iX$; the $\iY$
process needs only the MLR property. Confirm this asymmetry is
intended.}

\begin{theorem}[Value of LR-better information]\label{thm:main}
Suppose
\begin{enumerate}
\item $r(\theta,a)$ is increasing in $\theta$ for all $a$ and has
  increasing differences,
\item $L_\iX(x|\theta,a)$, $L_\iY(y|\theta,a)$, and
  $g(\theta_{k-1}|\theta_k,a)$ satisfy the MLR property for all $a$,
\item $L_\iX(x|\theta,a)$ and $g(\theta_{k-1}|\theta_k,a)$ satisfy the
  stochastic version of increasing differences, and
\item $\iX \LRI \iY$.
\end{enumerate}
Then, for all $k$ and all priors $\pi$,
\begin{equation}
V_{\iX,k}^*(\pi) \geq V_{\iY,k}^*(\pi) \ . \label{eq:main_result}
\end{equation}
\end{theorem}

\begin{proof}{Proof.}
We proceed by induction on $k$. For $k=0$,
$V_{\iX,0}^*(\pi) = V_{\iY,0}^*(\pi) = 0$. Now suppose
$V_{\iX,k-1}^*(\pi) \geq V_{\iY,k-1}^*(\pi)$ for all $\pi$. Fix a
prior $\pi$ and an action $a$. By the induction hypothesis,
\begin{equation}
\ev{V_{\iY,k-1}^*(\Pi_\iY(\pi,\tilde{y},a))} \leq
\ev{V_{\iX,k-1}^*(\Pi_\iY(\pi,\tilde{y},a))} \ .
\label{eq:induction_step}
\end{equation}
Next, take $u(\pi',a') = V_{\iX,k-1}(\pi',a')$ in
Lemma~\ref{lemma:LRI-ContVal}; this $u$ has LR-increasing differences
by Proposition~\ref{prop:monotone_policies}(i), whose hypotheses hold
for $\iX$ by assumptions (i)--(iii), and
$u^*(\pi') = V_{\iX,k-1}^*(\pi')$. The lemma then gives
\begin{equation}
\ev{V_{\iX,k-1}^*(\Pi_\iY(\pi,\tilde{y},a))} \leq
\ev{V_{\iX,k-1}^*(\Pi_\iX(\pi,\tilde{x},a))} \ .
\label{eq:lemma_step}
\end{equation}
Combining \eqref{eq:induction_step} and \eqref{eq:lemma_step} with the
recursion \eqref{eq:Vka}, we have, for every action $a$,
\begin{align}
V_{\iY,k}(\pi,a)
&= r(\pi,a) + \delta\, \ev{V_{\iY,k-1}^*(\Pi_\iY(\pi,\tilde{y},a))}
  \notag\\
&\leq r(\pi,a) + \delta\, \ev{V_{\iX,k-1}^*(\Pi_\iX(\pi,\tilde{x},a))}
  = V_{\iX,k}(\pi,a) \ .
\end{align}
Maximizing over $a \in A$ yields
$V_{\iY,k}^*(\pi) \leq V_{\iX,k}^*(\pi)$, completing the induction.
\Halmos
\end{proof}

Theorem~\ref{thm:main} is the monotone-POMDP counterpart of the
classical Blackwell comparison: it trades the universal quantifier over
decision problems for a larger set of comparable signal processes.
Because a Blackwell garbling is a special case of an LR-better garbling
(with a state-independent kernel), Theorem~\ref{thm:main} contains, for
monotone POMDPs, the familiar result that garbling in the sense of
Blackwell lowers the value function.

% ---------------------------------------------------------------------
% Section 4: Relationships among information orders.
% Drafted by Claude from Kim (2023, JET), per the authors' instruction,
% with \Claude{} comments flagging items for the authors to verify.
% Kim's definitions and results are restated in this paper's notation:
% states theta, signal processes X (better) and Y (worse); Kim's F and
% G correspond to our L_X and L_Y. Static comparison: the action/period
% arguments are suppressed throughout this section.
% ---------------------------------------------------------------------
\section{Relationships Among Information Orders}\label{sec:orders}

In this section, we relate the LR-better order to three information
orders from the literature: Blackwell's garbling order
\citep{blackwell1951comparison}, Lehmann's accuracy order
\citep{lehmann1988comparing}, and the monotone quasi-garbling order of
\citet{kim2023comparing}. Throughout this section, we compare the
signal processes in a single period and suppress the action argument;
in the dynamic model, the comparisons are assumed to hold for each
action. We write $F_\iX(x|\theta)$ for the cumulative distribution
function associated with the likelihood $L_\iX(x|\theta)$, and
similarly for $\iY$.

We begin by recalling the three orders. \citet{kim2023comparing}
introduces the monotone quasi-garbling order by relaxing Blackwell's
requirement that the added noise be independent of the state, while
restricting the state dependence to be \emph{reversely monotone} in
the sense of FOSD.

\begin{definition}[Garbling and monotone quasi-garbling]\label{def:BW-MQG}\quad
\begin{enumerate}
\item \citep{blackwell1951comparison} $\iX$ is \emph{Blackwell-better}
  than $\iY$ ($\iX \BW \iY$) if there exists a conditional probability
  distribution $P(y|x)$ such that
  $L_\iY(y|\theta) = \int P(y|x) L_\iX(x|\theta) \dif x$ for all
  $y \in Y$ and $\theta$.
\item \citep[Definition~3]{kim2023comparing} $\iX$ is a \emph{monotone
  quasi-garbling improvement} over $\iY$ ($\iX \MQG \iY$) if there
  exists a conditional probability distribution $P(y|x,\theta)$ such
  that (a) $L_\iY(y|\theta) = \int P(y|x,\theta) L_\iX(x|\theta)
  \dif x$ for all $y \in Y$ and $\theta$, and (b)
  $P(\,\cdot\,|x,\theta_1) \FOSD P(\,\cdot\,|x,\theta_2)$ for all
  $x \in X$ and $\theta_2 \geq \theta_1$.
\end{enumerate}
\end{definition}

\begin{definition}[Lehmann accuracy]\label{def:Lehmann}
\citep{lehmann1988comparing, kim2023comparing} $\iX$ is \emph{more
Lehmann accurate} than $\iY$ ($\iX \Leh \iY$) if, for all $y \in Y$,
the function
$\phi(\theta; y) = \sup\{x \,|\, F_\iX(x|\theta) \leq
F_\iY(y|\theta)\}$ is increasing in $\theta$.
\end{definition}

Lehmann's order has a simple interpretation: $\phi(\,\cdot\,; y)$ maps
a signal $y$ from $\iY$ to the signal from $\iX$ at the same
conditional quantile; $\iX$ is more accurate if this
quantile-matching map moves in the same direction as the state.
\citet{lehmann1988comparing} shows that, when payoffs and likelihoods
have appropriate monotone structure, more accurate experiments are
more valuable in single-decision problems; \citet{persico2000information}
and \citet{quah2009comparative} extend the reach of this result to
broader classes of monotone decision problems.
\citet{kim2023comparing} shows that, under the MLR property, Lehmann's
order is \emph{equivalent} to the monotone quasi-garbling order, which
gives Lehmann's condition a garbling interpretation.

The definitions of the four orders differ only in what they require of
the garbling kernel $P$ (or, for Lehmann's order, in avoiding a kernel
altogether): Blackwell requires state independence, the LR-better
order requires the state dependence to be reversely monotone in the LR
sense, and monotone quasi-garbling requires it to be reversely
monotone in the weaker FOSD sense. The following proposition collects
the resulting relationships.

\begin{proposition}[Relationships among the orders]\label{prop:orders}\quad
\begin{enumerate}
\item $\iX \BW \iY$ implies $\iX \LRI \iY$, and $\iX \LRI \iY$ implies
  $\iX \MQG \iY$. Neither implication can be reversed.
\item If $L_\iX$ satisfies the MLR property, then $\iX \MQG \iY$ if
  and only if $\iX \Leh \iY$ \citep[Theorem~1]{kim2023comparing}.
\item Without the MLR property on $L_\iX$, neither $\iX \BW \iY$ nor
  $\iX \MQG \iY$ implies $\iX \Leh \iY$
  \citep[Proposition~4]{kim2023comparing}.
\end{enumerate}
\end{proposition}

\begin{proof}{Proof of (i).}
If $\iX \BW \iY$ with a state-independent kernel $P(y|x)$, then
$P(y|x,\theta) = P(y|x)$ satisfies condition (ii) of
Definition~\ref{def:LRI} with equality, so $\iX \LRI \iY$. If
$\iX \LRI \iY$, then condition (ii) of Definition~\ref{def:LRI} states
that $P(\,\cdot\,|x,\theta_1) \LR P(\,\cdot\,|x,\theta_2)$ for
$\theta_2 \geq \theta_1$; since LR dominance implies FOSD, the same
kernel satisfies Definition~\ref{def:BW-MQG}(ii), so $\iX \MQG \iY$.
For the strictness of the first implication, the Normal example of
Section~\ref{sec:examples} with a strictly decreasing bias function
$d$ provides an LR-better garbling that is not a Blackwell garbling.
\Claude{Strictness claims to be completed: (a) for
$\BW \Rightarrow \LRI$, we should verify formally that the Normal
example with decreasing $d$ is not Blackwell-comparable (intuitively
the noise is state dependent, but a different state-independent kernel
could conceivably generate the same $L_\iY$---this needs an argument);
(b) for $\LRI \Rightarrow \MQG$, we need an example of a garbling
kernel that is reversely FOSD-monotone but not reversely
LR-monotone---and, ideally, such that no LR-monotone kernel generates
the same $\iY$. Your note at the end of the Uniform example in the
source file (mixtures of deterministic garblings, and uniform
garblings on $[x, x+d(\theta)]$) suggests candidate constructions.
Please advise which example you would like to develop.}
\Halmos
\end{proof}

Under the MLR property, then, the orders are nested:
\begin{equation}
\iX \BW \iY \ \Rightarrow\ \iX \LRI \iY \ \Rightarrow\
\iX \MQG \iY \ \Leftrightarrow\ \iX \Leh \iY \ .
\label{eq:nesting}
\end{equation}
\citet{kim2023comparing} shows that, in \emph{static} monotone
decision problems, the monotone quasi-garbling order characterizes the
value of information: it is sufficient for one signal process to yield
a higher expected payoff for all decision makers with monotone
decision problems (his Theorem~2), and it is also necessary when the
class of decision problems is rich enough to include all simple
hypothesis-testing problems (his Theorem~3 and Corollary~1). In view
of \eqref{eq:nesting}, one might hope that the same order would
suffice for \emph{dynamic} monotone decision problems. The reason we
work with the stronger LR-better order is the same reason the LR order
replaces FOSD throughout Section~\ref{sec:model}: an FOSD comparison
between garbling distributions need not survive Bayesian updating and
state transitions, while the LR comparison does. The two-state
example of Section~\ref{sec:tech-example} illustrates what can go
wrong when the reversely monotone structure fails; whether the
monotone quasi-garbling order itself is sufficient for value
comparisons in monotone POMDPs---or whether the gap between
$\LRI$ and $\MQG$ is essential in dynamic problems---is an open
question we intend to settle. \Claude{This last sentence states the
research question implied by your outline; if you already have a
counterexample showing MQG is \emph{not} sufficient in the dynamic
model, it belongs here and would sharpen the paper's contribution
considerably. Please advise.}

Two special cases deserve mention. First, when there are only two
states, every pair of MLR signal processes ordered by any of the above
criteria is comparable in a particularly transparent way;
\citet{jewitt2007information} shows that, under the MLR property,
Lehmann's order is equivalent to Blackwell dominance holding on all
dichotomies (all two-state restrictions) of the model. Second, your
outline records the claim that in models with \emph{two actions}, the
four orders coincide under the MLR property.
\Claude{I have not found this two-action equivalence stated in Kim
(2023)---his results connect MQG, Lehmann, and Blackwell but not in a
form specific to binary-action problems; the dichotomy result of
Jewitt (2007) concerns two \emph{states}, not two actions. If the
two-action claim is a result you intend to prove, I will set it up as
a proposition with a proof to be supplied; please clarify the intended
statement (e.g., ``for two-action monotone decision problems, $\iX$
yields a higher value than $\iY$ for all priors and payoffs if and
only if $\iX \MQG \iY$'') and its provenance.}

% ---------------------------------------------------------------------
% Section 5: Applications and examples.
% Sources: LRMQG.tex table of POMDP examples (inserted as-is per the
% authors' instruction, including the authors' margin notes); Normal/
% Poisson/Uniform garbling constructions; two-state technology adoption
% example with primitives recovered from 2statesTechAdoptionClaude.xlsx.
% ---------------------------------------------------------------------
\section{Applications}\label{sec:applications}

\subsection{Three Monotone POMDPs}\label{sec:three-models}

Table~\ref{tab:pomdp-examples} illustrates the scope of the monotone
POMDP framework with three models from the literature: the technology
adoption model of \citet{ulu2009uncertainty}, the machine replacement
model of \citet{white1979optimal} (see also \citealt{lovejoy1987some}),
and the prostate biopsy referral model of \citet{zhang2012optimization}.
In each model, the state, actions, and signals are ordered, the value
function is monotone in the beliefs, and the optimal policy is
monotone: as beliefs about the technology improve, the DM moves toward
adoption; as beliefs about the machine's condition worsen, the DM moves
toward replacement; and as beliefs point toward cancer, the DM moves
toward biopsy. Theorem~\ref{thm:main} applies to each model: improving
the signal process in the LR-better sense increases the optimal value
function. \Claude{Transitional text drafted; the table is inserted
as-is from your notes, including the two ``?'' cells and your margin
notes, per your instruction.}

\begin{landscape}
\begin{table}[h!]
\footnotesize
\centering
\caption{POMDP examples: \citet{ulu2009uncertainty},
\citet{white1979optimal}, \citet{zhang2012optimization}.}
\label{tab:pomdp-examples}
\renewcommand{\arraystretch}{1.4}
\begin{tabular}{>{\bfseries}p{2.5cm} p{6.5cm} p{6.5cm} p{5.5cm}}
\toprule
 & \textbf{Technology Adoption} & \textbf{Machine Replacement} & \textbf{Prostate Biopsy} \\
State of the world
    & $\theta_k \in \Theta \subseteq \Re, $ \newline technology benefit in period $k$
    & $s_t \in S=\{0,\ldots,n=\text{good as new}\}$, \newline machine state in period $t$
    & $s_t\in \{NC, C, T, M, D\}$, $NC$: no cancer, $C$: undetected prostate cancer,
    $T$: Treated, nonmetastatic (obs.), \newline $M$: metastasis (obs.), $D$: Death (obs.)\\
\midrule
Prior
  &$\pi(\theta_k)$, \newline continuous probability distn.
  & $\pi(s_t)$, \newline discrete probability distn.
  & $\pi_t \in [0,1]$, \newline probability of undetected cancer \\
\midrule
Signal process
    & $L(x|\theta_k)$ probability of observing signal $x$ given technology state $\theta_k$ \newline
    satisfies the MLR property
    &$r_{jk}^a$: probability of observing $k\in O$ given state $s_{t+1}=j$ and action $a_t=a$
  & $q_t(l_t|s_t)$: probability of observing PSA result $l_t \in \{1,\ldots,m\}$
    given state $s_t$ \\
\midrule
Actions $(a_k)$
  & $a_k = $ technology choice in period\newline
    $\in \{$adopt, investigate, reject$\}$\newline
    where adopt $\succ$ investigate $\succ$ reject
  & $a_t \in \{a:$ produce, $a^\prime:$ replace$\}$\newline
  where  $a^\prime\succ a$
  & $a_t = $ biopsy decision in period  $\in \{B, W\}$\newline
   where $W\succ B$\\
\midrule
Rewards\newline$(r_k(\theta_k, a_k))$
  & $r_k(\theta_k, a_k)$\newline
    $= \begin{cases} A\theta_k - K & \text{if } a_k = \text{adopt} \\ -c & \text{if } a_k = \text{investigate} \\ 0 & \text{if } a_k = \text{reject} \end{cases}$
  & $g(i,a)$: reward from action $a$ in state $s_i$, \newline
    has isotone differences
  & $\bar{r}_t(s_t, a_t)$\newline
    $= \begin{cases} 1 & \text{if } a_t = W \\ 1-\mu & \text{if } a_t = B, s_t=NC \\ 1-\mu-f_\epsilon & \text{if } a_t = B, s_t=C\end{cases}$ \newline{\CU{There are more states.}}\\
\midrule
Expected Rewards\newline$(r_k(\pi, a_k))$
  & $r_k(\pi, a_k)=\int_{\theta \in \Theta} r_k(\theta,a_k) \pi(\theta) \dif \theta$
  & $\sum_{i\in S} \pi_i g(i,a)$
  & $r_t(\pi_t, a_t)$\newline
    $= \begin{cases} 1 & \text{if } a_t = W \\ R_t(\pi_t) & \text{if } a_t = B\end{cases}$\\
\midrule
Transitions\newline$(\Pi(\theta_k;\pi, x,a_k))$
  & $\Pi(\theta_k;\pi, x,a_k)$\newline
    $= \begin{cases} ? & \text{if } a_k = \text{adopt or reject} \\ \frac{ L(x|\theta,a) \eta(\theta; \pi, a)}{f(x;\pi, a)}& \text{otherwise} \end{cases}$\newline
    where $\eta(\theta_k; \pi, a)=\int_{\theta_k \in \Theta} g(\theta_{k-1}|\theta_k, a) \pi(\theta_k) \dif \theta_k$ and $f(x;\pi, a)= \int_\theta L(x|\theta,a) \eta(\theta; \pi, a)\,\dif \theta$
  & $T(\pi, a, k)$\newline
   $=\begin{cases} ? & \text{if } a_k = \text{produce} \\ e_n \footnotesize{\text{(a unit vector with a 1 in the $n^{th}$ position)}}& \text{if } a_k=\text{replace} \end{cases}$
  & $\pi_{t+1}(s_{t+1})=\frac{q_{t+1}(l_{t+1}|s_{t+1})\sum_{s_t \in S}p_t(s_{t+1}|s_t, a_t)\pi_t(s_t)}{\sum_{s_{t+1}\in S}q_{t+1}(l_{t+1}|s_{t+1})\sum_{s_t \in S}p_t(s_{t+1}|s_t, a_t)\pi_t(s_t)}$\newline \CU{we can use something like our notation here.}\\
\midrule
Properties of the value function
  & $V_k^*(\pi)$ is $LR$-increasing and convex in $\pi$, \newline
    $V_k(\pi, a_k)$ satisfies $LR$ increasing differences.
  & For each $s_k$, $v(s_k, p_k; \sigma)$ is increasing\newline
    and convex in $p_k$ and increasing in\newline
    $\sigma$
  & $v_t(\pi_t)$ is decreasing and convex in $\pi_t$,\newline $R_t(\pi_t)$ linear and decreasing in $\pi_t$, \newline Control-limit type policy optimal \\
\bottomrule
\end{tabular}
\end{table}
\end{landscape}

\subsection{Constructing LR-Better Signal Processes}\label{sec:examples}

We next present three families of examples in which an LR-better
garbling arises naturally: a Normal signal with a state-dependent
bias, a Poisson signal with state-dependent thinning, and a Uniform
signal with state-dependent scaling. In each case, the garbled process
would typically \emph{not} be comparable to the original in
Blackwell's order because the added noise depends on the state, yet
the LR-better order applies and, by Theorem~\ref{thm:main}, a monotone
POMDP has a higher value function with the original signal process.

\subsubsection{Normal Likelihood with a State-Dependent Bias.}
\label{sec:normal-example}
Let $F(x|\mu, \sigma^2)$ be a Normal signal-generating process with
mean $\mu$ and variance $\sigma^2$, where the state of the world is
the unknown mean $\mu$. This likelihood satisfies the MLR property.
Suppose that we add a bias term that depends on $\mu$ and observe
\begin{equation}
Y = X + d(\mu) + Z \label{eq:normal-garbling}
\end{equation}
instead, where $d(\mu)$ is a decreasing function of $\mu$ and $Z$ is a
Normally distributed error term with mean $0$ and variance $\nu^2$,
independent of $X$. The garbling kernel is then
\begin{equation}
H(y|x, \mu) \sim N(x+d(\mu),\, \nu^2) \ ,
\end{equation}
a conditional probability distribution that is LR-decreasing in $\mu$
because $d(\mu)$ is decreasing; for example, $d(\mu) =
-\delta\mu/\sigma$ for some $\delta > 0$. Given the definition of the
biased signal process $\iY$, we have
\begin{equation}
G(y|\mu) \sim N(\mu+d(\mu),\, \sigma^2+\nu^2) \ ,
\end{equation}
which also satisfies the MLR property provided that $\mu + d(\mu)$ is
increasing; with $d(\mu) = -\delta\mu/\sigma$, this holds if
$\delta \leq \sigma$. Then $\iX \LRI \iY$: since both likelihoods
satisfy the MLR property, Theorem~\ref{thm:main} implies that a
monotone POMDP has a higher value function with likelihood $F$ than
with $G$. If $d(\mu)$ is constant, the noise is state independent and
$F$ Blackwell-dominates $G$; a strictly decreasing $d$ takes the
comparison outside Blackwell's order but keeps it within the LR-better
order.

The bias term has a natural interpretation: signals about better
technologies may be systematically dampened---for example, reviews of
genuinely superior products may understate their advantage---and such
skepticism-type distortions degrade information in exactly the
reversely monotone way the LR-better order captures.
\Claude{Interpretation sentence drafted by me---replace or delete if
it does not match the story you want to tell.}

\subsubsection{Poisson Likelihood with State-Dependent Thinning.}
\label{sec:poisson-example}
Let $F(x|\lambda)$ be a Poisson distribution with rate $\lambda$,
which has the MLR property in $\lambda$. Suppose that, instead of
observing the count $X$, each arrival is observed independently with
probability $d(\lambda)$. The garbling distribution $H(y|x,\lambda)$
is then Binomial with $x$ trials and success probability
$d(\lambda)$; it is LR-decreasing in $\lambda$ if $d(\lambda)$ is
decreasing. We now show that, after this thinning, the resulting
signal distribution is Poisson with rate $d(\lambda)\lambda$:
\begin{align}
g(y|\lambda) &= \sum_{x=y}^\infty h(y|x, \lambda)\, f(x|\lambda) \\
&= \sum_{x=y}^\infty \binom{x}{y}\, [d(\lambda)]^y
  [1-d(\lambda)]^{x-y}\, \frac{e^{-\lambda}\lambda^x}{x!} \\
&= \frac{e^{-\lambda}[d(\lambda)]^y}{y!} \sum_{x=y}^\infty
  \frac{\lambda^x [1-d(\lambda)]^{x-y}}{(x-y)!} \\
&= \frac{e^{-\lambda}[d(\lambda)]^y}{y!} \sum_{z=0}^\infty
  \frac{\lambda^{z+y} [1-d(\lambda)]^{z}}{z!} \\
&= \frac{e^{-\lambda}[d(\lambda)\lambda]^y}{y!} \sum_{z=0}^\infty
  \frac{\left(\lambda[1-d(\lambda)]\right)^{z}}{z!} \\
&= \frac{e^{-\lambda}[d(\lambda)\lambda]^y\,
  e^{\lambda(1-d(\lambda))}}{y!} \\
&= \frac{e^{-\lambda d(\lambda)}[d(\lambda)\lambda]^y}{y!} \ .
\end{align}
The second equality uses the definitions of the Binomial and Poisson
distributions; the third uses $\binom{x}{y} = \frac{x!}{(x-y)!\,y!}$;
the fourth substitutes $z = x-y$; the fifth collects the powers of
$\lambda$; and the sixth uses $\sum_{z=0}^\infty x^z/z! = e^x$. The
resulting Poisson distribution with rate $d(\lambda)\lambda$ has the
MLR property if $d(\lambda)\lambda$ is increasing in $\lambda$. Thus,
if $d$ is decreasing but $d(\lambda)\lambda$ is increasing, we have
$\iX \LRI \iY$ with both processes satisfying the MLR property, and
Theorem~\ref{thm:main} applies. This is a model of state-dependent
undercounting: arrivals from better processes (higher $\lambda$) are
individually less likely to be recorded.

\subsubsection{Uniform Likelihood with State-Dependent Scaling.}
\label{sec:uniform-example}
Let $F(x|\theta)$ be a Uniform distribution on $[0, \theta]$, which
has the MLR property in $\theta$. Suppose that instead of observing
$X$ we observe
\begin{equation}
Y = d(\theta) X \ .
\end{equation}
Then $Y|\theta$ is Uniform on $[0, d(\theta)\theta]$, which has the
MLR property if $d(\theta)\theta$ is increasing in $\theta$. For this
transformation, the garbling distribution is deterministic:
\begin{equation}
H(y|x,\theta) = \left\{
\begin{array}{cl}
1 & \mbox{if } y = x\, d(\theta)\ , \\
0 & \mbox{otherwise}\ .
\end{array}\right.
\end{equation}
To see that $H(y|x,\theta)$ is LR-decreasing when $d(\theta)$ is
decreasing in $\theta$, consider the inequality
\begin{equation}
H(y^2|x,\theta^2)\, H(y^1|x,\theta^1) \leq
H(y^2|x,\theta^1)\, H(y^1|x,\theta^2) \label{eq:uniform-LR}
\end{equation}
for $y^1 \leq y^2$ and $\theta^1 \leq \theta^2$. If
$y^1 = x\, d(\theta^2) \leq y^2 = x\, d(\theta^1)$, the inequality
holds because the left side is $0$ and the right side is $1$; for any
other selection of $y^1 \leq y^2$, both sides equal zero.
\CU{I think the above example would work for $Y = X + d(\theta)$. I
tried more general garbling distributions $H(y|x,\theta)$ that are not
deterministic. For example, we can observe $x$ or $x\,d(\theta)$ with
equal probability (does not satisfy the LR-decreasing property), or
$H(y|x,\theta)$ uniform between $x$ and $x+d(\theta)$ (is
LR-decreasing if $d(\theta)$ is decreasing, but the marginal for
$y|\theta$ is a trapezoidal distribution and it is not clear if it can
be made to have the MLR property).}

\subsection{A Two-State Technology Adoption Example: When
State-Increasing Garbling Yields No Ranking}\label{sec:tech-example}

The reversely monotone structure of the LR-better order---noise that
degrades the signal more in \emph{higher} states---is not a technical
artifact: when the state dependence of the noise runs in the opposite
direction, the two signal processes may genuinely be incomparable. We
illustrate this with a simple two-state technology adoption example.

The technology's benefit is $\theta \in \{0.1, 0.7\}$; we refer to
these as the low and high states. The DM's prior probability on the
high state is $0.4$. The DM either adopts the technology, earning
$A\theta - K$ with $A = 1$ and adoption cost $K = 0.3$ (i.e., $-0.2$
in the low state and $0.4$ in the high state), or rejects it, earning
$0$. Before deciding, the DM observes the outcome of two independent
trials of the technology, each of which succeeds with probability
$\theta$.

Under signal process $\iX$ (experiment $F$), the DM observes the
number of successes $x \in \{0,1,2\}$, so $L_\iX(x|\theta)$ is
Binomial with two trials and success probability $\theta$; this
likelihood satisfies the MLR property. Under signal process $\iY$
(experiment $G$), each success is \emph{recorded} only with
probability $d(\theta)$, with $d(0.1) = 0.4$ and $d(0.7) = 0.7$: the
probability that a success is observed \emph{increases} with the
state. The DM observes the number of recorded successes
$y \in \{0,1,2\}$, so $L_\iY(y|\theta)$ is Binomial with two trials
and success probability $\theta\, d(\theta)$ ($0.04$ in the low state
and $0.49$ in the high state), which also satisfies the MLR property.
The garbling kernel taking $x$ to $y$ is Binomial with $x$ trials and
success probability $d(\theta)$---exactly as in the Poisson thinning
example of Section~\ref{sec:poisson-example}, except that here $d$ is
\emph{increasing} in the state, so the kernel is LR-\emph{increasing}
in $\theta$ and the LR-better condition (Definition~\ref{def:LRI}(ii))
fails; indeed, even the weaker FOSD condition of monotone
quasi-garbling fails.

Figure~\ref{fig:tech-trees} shows the decision trees for the two
experiments. At the prior of $0.4$, experiment $F$ is more valuable:
the DM's expected value is $0.123$ under $F$ versus $0.109$ under $G$.
But this ranking does not hold for all priors: for sufficiently
pessimistic priors, $G$ is strictly more valuable (for example, at a
prior of $0.01$ on the high state, the DM's value is $0.0006$ under
$G$ and $0$ under $F$---under $F$, no signal is favorable enough to
justify adoption, while under $G$ the strongest signal is). The two
value functions cross as a function of the prior, so neither
experiment dominates the other. \Claude{Values computed from your
spreadsheet (2statesTechAdoptionClaude.xlsx): the value-by-prior table
shows $G$ strictly better for priors around $0.01$ and $F$ strictly
better for moderate priors, e.g., $0.1228$ vs.\ $0.10898$ at prior
$0.4$.}

Why does no ranking exist? When the observation probability increases
with the state, the signal under the signal process $G$ carries two
pieces of evidence rather than one: a recorded success indicates not
only that the trial succeeded, but also that the state was one in
which successes are more likely to be recorded. These two pieces of
evidence point in the same direction and reinforce each other, so a
favorable signal under $G$ is stronger evidence for the high state
than the corresponding signal under $F$. The price is paid at the
other end of the signal scale: silence becomes ambiguous, since it may
reflect either a failed trial or an unobserved success, and therefore
constitutes weaker evidence for the low state than it does under $F$.
The effect on beliefs is a one-sided stretch: posteriors move farther
than under $F$ after good news, but not as far after bad news. $G$ is
thus genuinely more informative for decisions that hinge on confirming
the high state and genuinely less informative for decisions that hinge
on detecting the low state, and neither experiment can dominate the
other: which one is more valuable depends on the prior beliefs.

\begin{figure}[t]
\FIGURE{%
\begin{tabular}{c@{\qquad}c}
\resizebox{0.45\textwidth}{!}{%
\begin{tikzpicture}[x=1.15cm,y=1.15cm]
\node[chance] (root) at (0,-3.6125) {};
\node[nval] at (root.south) {$0.123$};
\payoffheader{0.75}
\sigblock{root}{0}{.522\ $x=0$}{.931}{.069}{$0$}{$-0.159$}{plain}{opt}
\sigblock{root}{-2.55}{.276\ $x=1$}{.391}{.609}{$0.165$}{$0.165$}{opt}{plain}
\sigblock{root}{-5.1}{.202\ $x=2$}{.030}{.970}{$0.382$}{$0.382$}{opt}{plain}
\end{tikzpicture}}
&
\resizebox{0.45\textwidth}{!}{%
\begin{tikzpicture}[x=1.15cm,y=1.15cm]
\node[chance] (root) at (0,-3.6125) {};
\node[nval] at (root.south) {$0.109$};
\payoffheader{0.75}
\sigblock{root}{0}{.657\ $y=0$}{.842}{.158}{$0$}{$-0.105$}{plain}{opt}
\sigblock{root}{-2.55}{.246\ $y=1$}{.187}{.813}{$0.288$}{$0.288$}{opt}{plain}
\sigblock{root}{-5.1}{.097\ $y=2$}{.010}{.990}{$0.394$}{$0.394$}{opt}{plain}
\end{tikzpicture}}
\\[4pt]
(a) Experiment $F$. & (b) Experiment $G$.
\end{tabular}}%
{Decision trees for the two-state technology adoption example
($\theta\in\{0.1,0.7\}$, prior on the high state $0.4$, adoption cost
$0.3$).\label{fig:tech-trees}}%
{Squares are decision nodes, circles are chance nodes, and bold
branches mark optimal decisions; node values appear below each node.}
\end{figure}

% ---------------------------------------------------------------------
% Section 6: Conclusion. Drafted by Claude; the outline's "Extensions:
% risk aversion" item is folded in as future work pending material.
% ---------------------------------------------------------------------
\section{Conclusion}\label{sec:conclusion}

Blackwell's comparison of experiments ranks information structures for
all decision problems, but ranks few pairs of information structures.
This paper develops a comparison tailored to monotone POMDPs---the
class of partially observed dynamic programs, common in operations
research, in which value functions and optimal policies are monotone
in likelihood-ratio-ordered beliefs. The LR-better order permits the
garbling that degrades a signal process to depend on the state,
requiring only that the garbling in a lower state likelihood-ratio
dominate the garbling in a higher state. The order is implied by
Blackwell sufficiency, implies the monotone quasi-garbling order that
characterizes the value of information in static monotone decision
problems \citep{kim2023comparing}, and---because likelihood-ratio
dominance survives Bayesian updating and state transitions---delivers
value-function comparisons at every horizon of a monotone POMDP. Our
examples show that the reversely monotone structure of the garbling is
essential: when noise instead intensifies with the state, value
functions can cross and no prior-free ranking exists.

Several extensions merit further study. First, our analysis assumes a
risk-neutral DM maximizing expected discounted rewards; extending the
comparison to risk-averse DMs would require tracking how the LR-better
order interacts with concave utility, paralleling the way
\citet{athey2018value} treat payoff functions with concave incremental
returns in static problems. \Claude{The outline lists ``risk
aversion'' as an extension but no material was provided; this sentence
is a placeholder framing---replace or expand when you decide the
scope.} Second, the exact gap between the LR-better and monotone
quasi-garbling orders in dynamic problems remains to be
characterized: is the stronger order necessary for value comparisons
in monotone POMDPs, or merely sufficient? Third, it would be useful to
develop further families of signal processes---beyond the Normal,
Poisson, and Uniform families presented here---for which the
LR-better comparison can be verified directly from model primitives.


% Acknowledgments to be added.

\bibliographystyle{informs2014}
\bibliography{refs}

\end{document}
